\section*{Hoofdstuk \ref{ch:g3:stones-pebbles-sand} --- Stenen, kiezels en zand}

\begin{solution}{exo:g3:stones-pebbles-sand:1}
Het \omterm{def:g3:stones-pebbles-sand:rock}{vaste gesteente} is het harde \omterm{def:g3:stones-pebbles-sand:rock}{gesteente} dat overal onder de bodem ligt.
We zien het waar geen bodem is om het te verbergen: op de steile flanken
van \omterm{def:g1:landscapes-around-us:landform}{bergen}, in kliffen aan \omterm{def:g2:water-on-earth:ocean}{zee}, langs wegen die door een \omterm{def:g1:landscapes-around-us:landform}{heuvel} zijn
uitgegraven.
\end{solution}

\begin{solution}{exo:g3:stones-pebbles-sand:2}
\omterm{def:g3:stones-pebbles-sand:sand}{Slib}, \omterm{def:g3:stones-pebbles-sand:sand}{zand}, \omterm{def:g3:stones-pebbles-sand:pebble}{kiezel}, rotsblok.
\end{solution}

\begin{solution}{exo:g3:stones-pebbles-sand:3}
1 millimeter is kleiner dan 2 millimeter: \omterm{def:g3:stones-pebbles-sand:sand}{zand}. 3 millimeter is groter
dan 2 millimeter: \omterm{def:g3:stones-pebbles-sand:pebble}{grind}.
\end{solution}

\begin{solution}{exo:g3:stones-pebbles-sand:4}
Natuurlijk: de \omterm{def:g3:stones-pebbles-sand:pebble}{kiezel}, het blok graniet (door mensen gezaagd, maar het
\omterm{def:g3:stones-pebbles-sand:rock}{gesteente} is natuurlijk), het \omterm{def:g3:stones-pebbles-sand:sand}{zand} op het strand. Door mensen gemaakt: de
baksteen en het glas.
\end{solution}

\begin{solution}{exo:g3:stones-pebbles-sand:5}
Het water van de \omterm{def:g1:landscapes-around-us:water}{rivier} heeft hem lang over de bodem voortgerold en tegen
andere stenen gestoten: zijn hoeken zijn afgesleten.
\end{solution}

\begin{solution}{exo:g3:stones-pebbles-sand:6}
Graniet, voor een muur; dunne, platte stenen platen (een \omterm{def:g3:stones-pebbles-sand:rock}{gesteente} zoals leisteen), voor daken.
\end{solution}

\begin{solution}{exo:g3:stones-pebbles-sand:7}
De ronde: hoe langer een steen door water wordt voortgerold, hoe meer zijn
hoeken afslijten. De scherpe heeft nog niet ver gereisd van de rots waar
hij van afbrak.
\end{solution}

\begin{solution}{exo:g3:stones-pebbles-sand:8}
De grootste \omterm{def:g3:stones-pebbles-sand:pebble}{kiezel} is 3 centimeter lang. Eén centimeter is 10 millimeter,
en $10 \div 2 = 5$: er passen vijf korrels van 2 millimeter op
1 centimeter.
\end{solution}

\begin{solution}{exo:g3:stones-pebbles-sand:9}
In de rustige vijver. Snel water houdt de piepkleine slibkorrels in
beweging; alleen waar het water bijna stilstaat, kunnen ze naar de bodem
zakken.
\end{solution}

\begin{solution}{exo:g3:stones-pebbles-sand:10}
Omdat ze gemaakt worden van materialen die uit de grond komen: bakstenen
van klei, beton van \omterm{def:g3:stones-pebbles-sand:sand}{zand}, \omterm{def:g3:stones-pebbles-sand:pebble}{grind} en cement, glas van \omterm{def:g3:stones-pebbles-sand:sand}{zand}. Klei, \omterm{def:g3:stones-pebbles-sand:sand}{zand} en
\omterm{def:g3:stones-pebbles-sand:pebble}{grind} komen van gebroken \omterm{def:g3:stones-pebbles-sand:rock}{gesteente}.
\end{solution}
